\nolabel{section}{Введение}

Современная разработка Android-приложений в парадигме чистой архитектуры \cite{refd} предполагает разделение системы на независимые слои: слой данных, слой домена и слой представления. Каждый слой оперирует собственными моделями данных. Переход между слоями требует преобразования одной модели в другую — для этого используются специальные функции-преобразователи, называемые мапперами. Мапперы могут принимать один или несколько параметров и возвращать целевую модель. В типичном проекте количество таких функций может достигать нескольких сотен, причём логика преобразования часто тривиальна и сводится к копированию значений одноимённых полей. Это порождает проблему шаблонного кода, которая увеличивает время разработки, снижает читаемость и повышает риск ошибок \cite{refi}.

Одним из эффективных способов борьбы с шаблонным кодом является \textbf{метапрограммирование} — подход, при котором программа пишет или модифицирует другую программу (или себя) в процессе компиляции или выполнения. Метапрограммирование позволяет автоматизировать повторяющиеся конструкции, генерировать код по описанию и адаптировать поведение под структуру данных без ручного дублирования. В экосистеме Kotlin \cite{refg} различают два вида метапрограммирования: \textbf{рефлексия} (анализ и вызов во время выполнения) и \textbf{генерация кода на этапе компиляции} (compile-time code generation). Рефлексия проста в использовании, но несёт накладные расходы и вызывает проблемы совместимости с обфускацией на Android. Генерация кода лишена этих недостатков и становится всё более популярной.

Наиболее современным инструментом для метапрограммирования в Kotlin является \textbf{Kotlin Symbol Processing (KSP)} \cite{refj} — библиотека от Google, работающая напрямую с Kotlin-символами. KSP позволяет создавать процессоры, которые анализируют аннотации и генерируют дополнительный исходный код на этапе компиляции, что идеально подходит для автоматического создания мапперов.
